\section{Introduction}
\label{Introduction}

Flooding can disrupt road access, increase evacuation demand, and constrain the use of evacuation centers. During a flood event, disaster planners and responders must determine which routes remain feasible, which centers can accommodate affected residents, and how available vehicles and other resources can be distributed. These decisions become more difficult when road conditions, shelter capacity, affected population, and resource availability change across possible flood situations.

Maps and geographic information can show locations, roads, and hazard areas, but static representations alone provide limited support for testing how changing conditions affect evacuation options. A scenario-based decision-support system can integrate relevant geographic and operational information so that planners can configure conditions, examine their consequences, and compare alternative evacuation plans before an event occurs.

This study develops \textit{OverFlow}, a geospatial decision-support system for flood-evacuation planning. The system is intended to help authorized local disaster planners and responders configure and compare scenarios involving rainfall, affected population, road or bridge closures, evacuation-center status and capacity, and available resources. It provides decision-support information on flood-affected areas, feasible and alternative evacuation routes, center allocation, estimated evacuation time, and resource allocation. The system supplements, rather than replaces, the judgment and authority of disaster-management personnel.

\subsection{Project Context}

Flooding can change the conditions under which an evacuation plan must operate. Roads and bridges may become unavailable, the number of people requiring assistance may vary between areas, evacuation centers may have limited capacity, and available vehicles and other resources may be insufficient for every possible situation. Local disaster planners must consider these conditions together when determining which routes remain usable, where affected residents can be accommodated, and what alternatives are available when an expected route or center cannot be used.

Maps can identify roads, boundaries, waterways, facilities, and hazard-related locations, but a map by itself does not show the consequences of changing several planning conditions at once. An evacuation plan may need to be reconsidered when rainfall conditions change, a road is manually closed, an evacuation center becomes unavailable or reaches capacity, the affected population increases, or fewer resources are available. Examining these conditions as configurable scenarios gives planners a structured way to compare possible evacuation arrangements before selecting or revising a plan.

This project develops \textit{OverFlow}, a localized two-dimensional simulation-based decision-support system for flood evacuation planning. The system is intended for authorized local disaster planners and responders. It combines terrain, road, boundary, evacuation-center, population, and resource information in an interactive planning environment. Users configure rainfall intensity and duration, affected population, road or bridge closures, evacuation-center status and capacity, and available resources. The system then presents flood-affected conditions and analyzes feasible and alternative routes, center allocation, estimated evacuation time, and resource distribution for comparison with other scenarios.

The flood component uses a digital elevation model as its terrain input. D8 flow direction and flow accumulation represent how terrain directs and concentrates water, while a simplified rainfall overlay and threshold classify areas for scenario analysis. Drainage infrastructure may be displayed as a reference layer but is not included in the flood computation. The rainfall-to-indicator formula, classification threshold, and validation procedure still require final definition; therefore, the output must be treated as a relative scenario representation rather than a calibrated hydraulic result or exact flood forecast.

OverFlow is intended to supplement the planning process by organizing relevant conditions and showing how changes in those conditions affect evacuation options. It does not make autonomous emergency decisions, guarantee that a route will remain safe during an actual event, or replace the authority and judgment of disaster-management personnel. The final locality, geographic coverage, operating organization, accessible datasets, routing and allocation methods, and evaluation design remain subject to confirmation.

\subsection{Purpose and Description}

The purpose of this study is to develop \textit{OverFlow}, a localized, scenario-based decision-support system for flood-evacuation planning. A planner configures rainfall intensity and duration, affected population by barangay, evacuation-center status, center capacity, manual road or bridge closures, and available resources. Using the geographic and terrain data for the selected study area, the system represents the resulting flood and evacuation conditions in an interactive two-dimensional planning environment.

For each scenario, the system identifies flood-affected areas and analyzes evacuation options under the configured constraints. It is intended to identify feasible and alternative routes, assign affected groups to available evacuation centers, estimate evacuation time, and present resource-allocation information. The planner can compare the results of different scenarios and use them as a reference in preparing or revising an evacuation plan.

The output is decision-support information, not an autonomous emergency decision or a guarantee of actual road safety. Results depend on the completeness, accuracy, resolution, and currency of the underlying data, as well as on the selected scenario inputs. The simplified terrain-based flood computation is suitable for comparing relative conditions, but it does not model detailed hydraulic behavior, flow velocity, infiltration, or drainage pipes and canals.

\subsection{Objectives of the Study}

The general objective of this study is to develop and evaluate a localized two-dimensional simulation-based decision-support system for creating and comparing flood-evacuation planning scenarios.

Specifically, it aims to:
\begin{enumerate}
  \item integrate validated local map, road, hazard, evacuation-center, population, and resource data for the selected study area;
  \item represent flood-affected areas from DEM-based terrain accumulation using rainfall intensity and duration, together with road disruption, evacuation demand, center capacity, and resource availability;
  \item analyze feasible evacuation routes, center allocation, estimated evacuation time, and resource distribution under configured scenario constraints;
  \item present scenario results through an interactive two-dimensional map and comparison interface; and
  \item evaluate the developed system using methods and metrics to be selected for the final study.
\end{enumerate}

\subsection{Significance of the Study}

The study is intended to provide an interactive way of examining localized flood-evacuation scenarios. By showing flood-affected areas together with road accessibility, evacuation demand, center capacity, and resources, \textit{OverFlow} may help users compare the consequences of alternative planning conditions before an emergency occurs.

For local disaster planners and responders, the system may provide a supplementary reference for preparing and revising evacuation plans. It can support the examination of feasible routes, alternate routes, center allocation, estimated evacuation time, and resource requirements under different scenario conditions. It does not replace existing disaster-management procedures or the judgment of authorized personnel.

For local engineering and planning personnel, barangay responders, and evacuation-center managers, the system may support coordination by presenting relevant geographic and capacity information in one planning environment. Community residents are intended indirect beneficiaries when better-informed preparedness and evacuation plans are developed. The study may also provide future researchers and developers with a reference for localized, terrain-based flood-evacuation decision-support systems.

\subsection{Scope and Delimitations of the Study}

This study covers the development of \textit{OverFlow}, a localized two-dimensional flood-evacuation planning and decision-support system. It includes configurable scenario inputs for rainfall intensity and duration, affected population by barangay, evacuation-center status and capacity, manual road or bridge closures, and available rescue resources. The system uses DEM data to compute terrain-based flood-affected areas through D8 flow direction, flow accumulation, rainfall overlay, and threshold-based classification. It represents roads, bridges, barangay boundaries, evacuation centers, population information, and resource conditions to support analysis and comparison of evacuation scenarios.

The system produces decision-support outputs that include flood-affected areas, recommended and alternative evacuation routes, center allocation, estimated evacuation time, resource-allocation information, dynamic maps, and scenario comparisons. Drainage infrastructure may be included only as a reference map layer. The final study area, deployment organization, and availability of local datasets remain subject to confirmation; the system's implemented coverage and features will be bounded by data that can be legitimately obtained and validated during the study.

The study excludes real-time flood forecasting, public navigation, emergency dispatch control, and direct control of physical infrastructure. It does not implement a calibrated hydraulic model and does not claim exact flood depth, velocity, or prediction of actual flood behavior. The model does not represent detailed hydraulic processes such as soil infiltration, drainage pipes and canals, or multi-directional flow spreading. Its outputs are intended for relative scenario comparison and human planning support, not as absolute emergency instructions.

\subsection{Operational Definition of Terms}

\textbf{Affected Population.} Refers to the number of people assigned to a barangay or area within a configured evacuation scenario.

\textbf{D8 Flow Direction.} Refers to a terrain-processing method that assigns each DEM cell a downhill flow direction toward its steepest neighboring cell among eight possible neighbors.

\textbf{Decision-Support System.} Refers to a system that helps a human decision-maker evaluate options. In \textit{OverFlow}, it presents processed information for examining flood-evacuation scenarios; it does not make autonomous emergency decisions.

\textbf{Digital Elevation Model (DEM).} Refers to a two-dimensional grid in which each cell stores ground elevation. It is the primary terrain input for the flood computation in \textit{OverFlow}.

\textbf{Drainage Infrastructure.} Refers to pipes, canals, and other structures that collect or convey water. In \textit{OverFlow}, available drainage information is a reference map layer and is not used to compute flooding.

\textbf{Evacuation-Center Allocation.} Refers to the assignment of affected groups to evacuation centers while considering accessibility and available capacity.

\textbf{Evacuation Route.} Refers to a road-network path considered usable between an affected area and an evacuation center under a scenario's constraints.

\textbf{Flood Scenario.} Refers to a configured set of conditions, including rainfall, affected population, road closures, center status and capacity, and available resources, used to examine a possible evacuation situation.

\textbf{Flow Accumulation.} Refers to a grid operation that counts or estimates the upstream area draining into a DEM cell, indicating locations where water may collect.

\textbf{OverFlow.} Refers to the proposed localized two-dimensional, simulation-based decision-support system for flood-evacuation planning developed in this study.

\textbf{Terrain-Based Flood Computation.} Refers to the simplified computation of flood-affected areas from rainfall input and DEM terrain data using D8 flow direction, flow accumulation, rainfall overlay, and threshold-based flood classification. It is not a hydraulic model.
